\scanpage{65}
$S$ schéma de car $p$.

$\mathcal C$ catégorie des schémas en groupes finis étales
\inserted{$E$} sur $S$ annulés par $p$
$\simeq$ catégorie des faisceaux \unsure{étales}
$\mathbb F_p$-vectoriels loc. libres de type fini.

$\mathcal C'$ catégorie des Modules loc. libres $M$ sur $S$, munis de
\[
F_M:M^{(p)}\xrightarrow{\sim}M.
\]
\[
\mathcal C\xrightarrow{\varphi}\mathcal C',
\qquad E\longmapsto
(E\otimes_{\mathbb F_p}\underline{\mathcal O}_S,
F_E\otimes_{\mathbb F_p}\underline{\mathcal O}_S)
\]
\[
\uline{\text{NB}}\quad
(E\otimes_{\mathbb F_p}\underline{\mathcal O}_S)^{(p)}
\simeq E^{(p)}\otimes_{\mathbb F_p}\underline{\mathcal O}_S
\xleftarrow[\unsure{f_{E/S}}]{\sim}
E\otimes_{\mathbb F_p}\underline{\mathcal O}_S)
\]

\uline{Proposition}\quad \uline{Le foncteur} $\varphi$
\uline{est une équivalence de catégories}.

a) \uline{Pleinement fidèle}. \deleted{\unsure{P}}
Par passage à des \uline{Hom}, revient à ceci:
\[
\Gamma(S,E)\xrightarrow{\sim}
\{x\in\Gamma(S,E\otimes_{\mathbb F_p}\mathcal O_S)
\mid F(x^{(p)})=x\}
\]
Par descente, on se ramène au cas $E$ constant, puis
$E=\mathbb F_{pS}$, et on est ramené alors à l'assertion que
\[
0-(\mathbb Z/p\mathbb Z)_S-G_{aS}
\xrightarrow{x\mapsto x^p-x}G_{aS}
\]
est exacte, ce qui est bien connu (il suffit de le voir sur
$\mathbb F_p$ par chgt de base \dots).

b) Essentiellement surjectif. Si $(M,F)$ est donné, il suffit
de voir qu'il est \uline{loc} de la forme voulue.
Considérons
\[
M\xrightarrow{\inserted{\unsure{\delta x}=F_Mf_Mx-x}}M
\]
on voit que c'est un \deleted{\unsure{iso}}morphisme étale de
schémas en groupes, donc son noyau est un schéma en groupes étale
$E$ sur $S$. Pour montrer qu'il est fini sur $S$, il suffit de
prouver qu'il est de \inserted{$\mathbb F_p$-}rang
\deleted{\unsure{égal}} \unsure{localt ct} \unsure{et un}
fini, égal au rang de $M$. Pour prouver ensuite que
$E\otimes_{\mathbb F_p}\mathcal O_S-M$ est un isom, il suffit de
le prouver fibre par fibre (géom). Bref, on est ramené au cas
\margin{75}
